\section*{Resumen}

En el contexto de la enseñanza del inglés como lengua extranjera, la clasificación de textos según su nivel de dificultad y la adaptación de materiales entre niveles constituyen dos necesidades centrales tanto para docentes como para estudiantes. Este trabajo, desarrollado en el marco del proyecto TaskGen, investiga y evalúa aproximaciones basadas en modelos de lenguaje de gran escala (LLMs) de código abierto para abordar ambas tareas según los niveles del Marco Común Europeo de Referencia para las Lenguas (CEFR).

Para la tarea de clasificación, se compararon múltiples LLMs de acceso abierto (familias LLaMA, Mistral y Gemma) y modelos propietarios (ChatGPT, DeepSeek) bajo 24 estrategias de prompting, que incluyen enfoques 0-shot, 1-shot y 2-shot, así como prompts enriquecidos con conocimiento de dominio (descriptores del CEFR y listas de vocabulario por nivel). Asimismo, se realizó un proceso de fine-tuning sobre el modelo Meta-Llama-3.1-8B-Instruct, y se propuso una métrica propia, la Accuracy Aproximada, que incorpora la naturaleza ordinal de la escala CEFR. Los resultados muestran que el modelo ajustado mediante fine-tuning supera de manera consistente a las estrategias basadas únicamente en prompting y a herramientas comerciales como Text Inspector, alcanzando un desempeño comparable al del clasificador de referencia de TSAR. Sin embargo, persiste una dificultad sistemática para distinguir niveles adyacentes, en particular en las fronteras B1--B2 y B2--C1.

Para la tarea de adaptación, enmarcada en la shared task TSAR 2025, se desarrolló el sistema TaskGen, basado exclusivamente en Llama-3.1-8B-Instruct, evaluando tres estrategias (runs): un refinamiento iterativo orientado a la preservación de contenido, un refinamiento orientado a la corrección de meta-información, y un sistema de \textit{ensemble} por votación heurística sobre diez prompts de distinta naturaleza. Esta última estrategia obtuvo el mejor desempeño global, posicionando a TaskGen en el 6º lugar del ranking general por equipos, destacándose especialmente en la preservación del significado (medida mediante MeaningBERT) tanto respecto al texto original como a las adaptaciones de referencia producidas por expertos.

En conjunto, este trabajo demuestra que los LLMs de código abierto de tamaño reducido (8B de parámetros), combinados con estrategias de prompting cuidadosamente diseñadas y recursos lingüísticos explícitos, constituyen una alternativa viable y competitiva frente a modelos privativos de mayor escala, lo cual resulta particularmente relevante para contextos de investigación con recursos computacionales limitados, como el entorno universitario argentino.